\section{问题五：绿电直连园区对电力系统的影响及政策建议}

\subsection{绿电直连园区容量渗透率提高的影响}

随着国家政策的鼓励推动\cite{ref1}，绿电直连园区将快速发展，其容量渗透率的显著提高将对所接入电力系统产生多方面影响\cite{ref7}。基于前四题的定量分析结论，从电力系统运行角度进行利弊分析。

\subsubsection{有利影响}

\textbf{（1）促进新能源就近消纳，降低弃风弃光率。}绿电直连模式使风光电力直接供给园区负荷，减少了远距离输电需求和通道建设投资。由问题三的计算结果可知，连续调节模式下园区自发自用比例可达60\%以上（产量63吨/日时全满足场景达20个），有效提升了新能源利用率。当大量园区采用该模式时，区域弃风弃光率将显著下降。

\textbf{（2）提供柔性负荷资源，缓解电网调峰压力。}制氢氨负荷具有功率连续可调的特性（10\%--100\%），相当于为电网提供了大规模的需求侧响应资源\cite{ref5}。由问题三的结果可知，单个园区的可调节容量达37.35~MW（$0.9 \times 41.5$~MW），当多个园区协同调节时，可为电网提供GW级的调峰能力，有效平抑风光出力波动对电网频率和电压的冲击。

\textbf{（3）推动氢能产业链发展，助力化工行业深度脱碳。}"绿电—绿氢—绿氨"一体化模式为化工行业提供了零碳原料路径\cite{ref4}，带动电解槽、储能等装备制造业发展。由问题四的分析可知，配置适当储能后园区可实现100\%风光利用率，生产的绿氨具有明确的碳减排价值。

\subsubsection{不利影响}

\textbf{（1）增加电网备用容量需求和尖峰负荷压力。}并网型园区在风光不足时段需从电网大量购电。由问题一的计算结果可知，单个园区的最大购电功率可达19.81~MW，多个园区同时购电将形成叠加尖峰负荷，增加系统备用容量需求和发电机组的启停成本。

\textbf{（2）加剧配电网电压波动和潮流反转风险。}园区在白天风光充裕时大量余电上网（最大售电功率36.59~MW），在夜间又大量购电，导致配电网潮流方向频繁反转，可能引起局部电压越限和保护装置误动作\cite{ref10,ref11}。渗透率越高，这一问题越突出。

\textbf{（3）增加电网规划和调度的不确定性。}绿电直连项目的选址、容量和运行模式多样化，且受风光资源随机性影响，使得配电网的负荷预测和网架规划难度增大。由问题二的24场景分析可知，不同风光条件下园区的购售电模式差异显著，给电网调度带来额外的不确定性。

\subsection{政策建议}

基于上述分析，提出以下推动绿电直连园区发展的政策建议：

\textbf{（1）完善绿电交易与碳交易联动机制。}建立绿电消纳量与碳配额的量化换算标准，使园区的绿电消纳行为能够获得碳市场收益，提升绿电直连项目的经济吸引力。由问题三的结果可知，连续调节模式下年均吨氨成本4283.62元/吨，若叠加碳交易收益可进一步降低至4000元/吨以下。

\textbf{（2）建立储能补贴和容量电价机制。}由问题四的分析可知，储能配置可将风光利用率从93.7\%提升至100\%，但储能投资成本较高。建议对配置储能的绿电直连项目给予投资补贴或容量电价补偿，降低储能配置的经济门槛。

\textbf{（3）制定差异化的绿电直连并网技术标准。}针对不同规模和类型的绿电直连项目，制定分级分类的并网技术标准，明确功率因数、谐波、电压偏差等技术指标要求，在保障电网安全的前提下简化并网审批流程。

\textbf{（4）推动园区与电网的协调调度机制。}建立园区柔性负荷参与电网调峰的市场化机制，使园区在响应电网调度指令时能够获得合理的经济补偿。由灵敏度分析可知，购电电价对吨氨成本的影响最大（灵敏度21.68\%），合理的峰谷电价信号能够有效引导园区主动参与电网调峰。

\subsection{基于 GRA-Entropy-TOPSIS 的园区多维度综合评价量化模型}

为客观评价电氢氨园区在不同模式和调节机制下的综合运行效益，量化系统隐性价值，本节构建灰色关联度分析（GRA）与熵权 TOPSIS 融合的多指标决策模型，对以下四种方案进行量化考量：
\begin{itemize}
    \item \textbf{方案A}：联网运行 + 离散调节（Q2 最优日产量方案）
    \item \textbf{方案B}：联网运行 + 连续调节（Q3 最优日产量方案）
    \item \textbf{方案C}：离网运行 + 无储能（Q4-1 最大出力方案）
    \item \textbf{方案D}：离网运行 + 最优储能配置（Q4-2 储能协同调度方案）
\end{itemize}

\subsubsection{评价指标体系设计}

模型选取以下 6 个涵盖经济性、合规性及系统隐性价值的维度构成评价指标体系：
\begin{enumerate}
    \item \textbf{吨氨综合成本 $X_1$（元/吨，成本型）}：反映方案的直接经济性。
    \item \textbf{全年绿电直连合规天数 $X_2$（天，效益型）}：衡量项目在政策合规层面的长期生命力。
    \item \textbf{风光自给消纳率 $X_3$（\%，效益型）}：反映园区对本地新能源的直接消纳能力。
    \item \textbf{电网支撑柔性调峰容量 $X_4$（MW，效益型）}：量化作为柔性负荷为电网提供的“削峰填谷”支撑价值。
    \item \textbf{碳减排强度与绿氨溢价 $X_5$（\%，效益型）}：反映双碳背景下碳减排强度及在未来的溢价潜力。
    \item \textbf{能源自治抗电价风险能力 $X_6$（\%，效益型）}：衡量在电网波动、电价上涨等外部风险下的能源安全系数。
\end{enumerate}

基于前文计算与合理估算，构建原始决策矩阵 $\mathbf{X}_{4 \times 6}$ 为：
\begin{equation}
\mathbf{X} = \begin{bmatrix}
4585.15 & 0.0 &   45.0 &  0.0  & 45.0 &  10.0 \\
4283.62 & 165.0 & 72.0 &  37.35 & 72.0 &  30.0 \\
4220.00 & 360.0 & 100.0 & 0.0  & 100.0 & 100.0 \\
4133.64 & 360.0 & 100.0 & 0.0  & 100.0 & 100.0 
\end{bmatrix}
\end{equation}

\subsubsection{数学建模步骤}

\textbf{步骤 1：决策矩阵标准化}。
在进行多指标决策分析时，由于各个评价维度的量纲与数量级差异显著（例如吨氨综合成本为数千元量级，而合规天数为数百天量级，自给消纳率等则为百分数），必须对其进行无量纲标准化处理，以消除不同量纲对决策结果造成的偏见。

定义标准化后的决策矩阵为 $\mathbf{Z} = (Z_{ij})_{m \times n}$。
\begin{itemize}
    \item 对于极大型（效益型）指标（如全年绿电直连合规天数、风光自给消纳率、调峰容量、碳减排与抗风险能力等），其规范化公式为：
    \begin{equation}
    Z_{ij} = \frac{X_{ij} - \min_{1 \le i \le m} X_{ij}}{\max_{1 \le i \le m} X_{ij} - \min_{1 \le i \le m} X_{ij}}
    \end{equation}
    \item 对于极小型（成本型）指标（如吨氨综合成本），其规范化公式为：
    \begin{equation}
    Z_{ij} = \frac{\max_{1 \le i \le m} X_{ij} - X_{ij}}{\max_{1 \le i \le m} X_{ij} - \min_{1 \le i \le m} X_{ij}}
    \end{equation}
\end{itemize}
通过此无量纲规范化处理，所有指标值均被线性映射到 $[0, 1]$ 区间内，保证了后续决策模型的严谨性与客观可信度。

\textbf{步骤 2：熵权法客观赋权}。
计算第 $j$ 个指标下第 $i$ 个方案的特征比重 $p_{ij} = \frac{Z_{ij}}{\sum_{i=1}^m Z_{ij}}$。各指标的熵值计算公式为：
\begin{equation}
e_j = -\frac{1}{\ln m} \sum_{i=1}^m p_{ij} \ln p_{ij}
\end{equation}
定义差异度 $d_j = 1 - e_j$，则各指标的客观权重为 $w_j = \frac{d_j}{\sum_{j=1}^n d_j}$。

\textbf{步骤 3：TOPSIS 相对贴近度计算}。
构建加权标准化决策矩阵 $\mathbf{V} = (V_{ij})_{m \times n} = (Z_{ij} \cdot w_j)_{m \times n}$。确定正理想解向量 $\mathbf{V}^+$ 和负理想解向量 $\mathbf{V}^-$：
\begin{equation}
V^+_j = \max_i V_{ij}, \quad V^-_j = \min_i V_{ij}
\end{equation}
计算各方案到正、负理想解的欧氏距离 $D^+_i = \sqrt{\sum_{j=1}^n (V_{ij} - V^+_j)^2}$，$D^-_i = \sqrt{\sum_{j=1}^n (V_{ij} - V^-_j)^2}$。则相对贴近度（TOPSIS得分）为：
\begin{equation}
S_i = \frac{D^-_i}{D^+_i + D^-_i}
\end{equation}

\textbf{步骤 4：灰色关联度（GRA）计算}。
以正理想解 $\mathbf{V}^+$ 为参考序列，计算第 $i$ 个方案在第 $j$ 个指标上的灰色关联系数：
\begin{equation}
\xi_{ij} = \frac{\min_i \min_j |V^+_j - V_{ij}| + \rho \max_i \max_j |V^+_j - V_{ij}|}{|V^+_j - V_{ij}| + \rho \max_i \max_j |V^+_j - V_{ij}|}
\end{equation}
取分辨系数 $\rho = 0.5$。加权灰色关联度为 $R_i = \sum_{j=1}^n w_j \xi_{ij}$。

\textbf{步骤 5：综合评分融合}。
结合 TOPSIS 在位置距离上的优势和 GRA 在形态相似度上的优势，计算综合评分：
\begin{equation}
Score_i = 0.5 S_i + 0.5 R_i
\end{equation}

\subsubsection{量化结果与深度剖析}

根据熵权法计算得到的权重分布如表\ref{tab:q5_weights}所示。

\begin{table}[H]
\centering
\caption{GRA-Entropy-TOPSIS 评价指标权重计算结果}\label{tab:q5_weights}
\begin{tabular}{llc}
\toprule
指标编号 & 指标名称 & 熵权权重(\%) \\
\midrule
$X_1$ & 吨氨综合成本 & 9.61 \\
$X_2$ & 全年绿电直连合规天数 & 10.90 \\
$X_3$ & 风光自给消纳率 & 10.64 \\
$X_4$ & 电网支撑柔性调峰容量 & 44.25 \\
$X_5$ & 碳减排强度与绿氨溢价 & 10.64 \\
$X_6$ & 能源自治抗电价风险能力 & 13.96 \\
\bottomrule
\end{tabular}
\end{table}

各方案的量化得分与排名如表\ref{tab:q5_ranking}所示。

\begin{table}[H]
\centering
\caption{各方案多维度综合量化评分与最终排名}\label{tab:q5_ranking}
\begin{tabular}{lccccc}
\toprule
评估方案 & TOPSIS得分 $S_i$ & GRA关联度 $R_i$ & 综合评分 $Score_i$ & 最终排名 \\
\midrule
\textbf{方案B（联网连续调节）} & 0.7543 & 0.8771 & 0.8157 & \textbf{1} \\
\textbf{方案D（离网最优储能）} & 0.3624 & 0.7050 & 0.5337 & \textbf{2} \\
\textbf{方案C（离网无储能）}   & 0.3562 & 0.6976 & 0.5269 & \textbf{3} \\
\textbf{方案A（联网离散调节）} & 0.0000 & 0.5168 & 0.2584 & \textbf{4} \\
\bottomrule
\end{tabular}
\end{table}

\textbf{量化结论分析}：
\begin{enumerate}
    \item \textbf{方案 B（联网连续调节）夺冠（综合评分 0.8157）}：方案 B 以绝对优势位居第一，根本原因在于其在「电网支撑柔性调峰容量」这一高权重维度（44.25\%）表现极其优异。连续调节使得园区能够作为 37.35~MW 级别的高价值需求侧响应负荷，主动参与配电网调峰。同时，方案 B 的吨氨成本（4283.62元/吨）仅比离网最低成本高出约 150 元/吨，但却提供了 36 吨/日的高稳定产量，实现了经济性与系统支撑价值的绝佳平衡。
    \item \textbf{离网方案（方案 D 与方案 C）位列中游}：方案 D（0.5337分）与方案 C（0.5269分）分列二、三位。虽然离网运行切断了与外部电网的联系，达成了 100\% 的风光自给消纳率、100\% 碳减排以及完美的抗电价风险能力（权重占比较高），且方案 D 配置储能后实现了 100\% 的新能源消纳与最低的吨氨成本（4133.64元/吨），但它们完全丧失了对公共电网的主动调峰支撑价值，高额的储能容量投资（410~MWh）也带来了巨大的资金占用，导致其在加权多维度评估中得分受限。
    \item \textbf{方案 A 垫底（0.2584分）}：离散调节模式下，由于设备开停机的刚性特征，不仅无法提供任何柔性调峰能力，还因频繁的功率突变对电网造成冲击。其极低的合规天数和最高的吨氨成本，导致其在各项指标上均处于绝对劣势。
\end{enumerate}

这一量化评估结果科学地揭示了未来绿电直连园区的演进方向：**联网运行配合连续柔性调节（方案 B）是当前技术与政策背景下的最优物理形态**，它能在最小化园区自身用能成本的同时，为新型电力系统提供极其珍贵的灵活性资源，实现园区与电网的双赢。

\section{模型评价与推广}

\subsection{模型优点}

（1）模型体系完整，从确定性分析（问题一）到离散优化（问题二）、连续优化（问题三）、再到含储能的两层优化（问题四），形成了递进式的建模框架，各层模型之间逻辑衔接紧密。

（2）求解方法高效且精确。问题二利用目标函数的可分解性，证明了贪心策略等价于整数线性规划的精确解，避免了组合爆炸问题。问题三将非线性问题转化为标准线性规划，保证了全局最优性。问题四采用两层优化结构，外层网格搜索保证了储能容量的全局搜索，内层线性规划保证了调度方案的最优性。

（3）模型具有明确的物理意义和工程可解释性。所有决策变量和约束条件均对应实际的物理量和工程约束，计算结果可直接指导园区的运行调度决策。

\subsection{模型不足}

（1）未考虑设备启停成本和最小运行时间约束。实际工程中，电解槽的频繁启停会加速催化剂老化，合成氨装置的启动需要预热时间，这些因素在本模型中未予考虑。

（2）风光出力的场景划分较为粗糙。24种场景（6风$\times$4光）难以完全覆盖实际风光资源的连续分布特征，可能遗漏极端场景。

（3）未考虑氢气储存环节。实际工程中制氢与合成氨之间通常设有氢气缓冲罐，允许两者在短时间尺度上解耦运行，这将为调度提供额外的灵活性。

\subsection{模型推广}

本模型框架可推广至以下场景：（1）其他类型的绿电直连项目，如绿电制甲醇、绿电炼钢等；（2）多园区协同优化，考虑多个园区共享电网接入点时的协调调度；（3）长时间尺度的投资规划问题，如风光装机容量和储能容量的联合优化配置。
